\section{Prepare a bootable micro-SD card}

The AM62x ROM code will look for boot images in a FAT partition on the
SD card. To be recognized by the ROM code, this partition must have a
special type and the {\em Bootable} flag set.

Let's prepare an SD card with such a partition.

Plug the SD card your instructor gave you on your workstation. Type
the \code{sudo dmesg} command to see which device is used by your
workstation. In case the device is \code{/dev/mmcblk0}, you will see
something like

\begin{verbatim}
[46939.425299] mmc0: new high speed SDHC card at address 0007
[46939.427947] mmcblk0: mmc0:0007 SD16G 14.5 GiB
\end{verbatim}

The device file name may be different (such as \code{/dev/sdb})
if the card reader is connected to a USB bus (either internally
or using a USB card reader).

In the following instructions, we will assume that your SD card is
seen as \code{/dev/mmcblk0} by your PC workstation.

Type the \code{mount} command to check your currently mounted
partitions. If SD partitions are mounted, unmount them:

\bashcmd{$ sudo umount /dev/mmcblk0p*}

We will erase the existing partition table by simply zero-ing the
first 16 MiB of the SD card:

\bashcmd{$ sudo dd if=/dev/zero of=/dev/mmcblk0 bs=1M count=16}

Now, let's use \code{cfdisk} to create the partitions that we need:

\begin{bashinput}
sudo apt install fdisk
sudo cfdisk /dev/mmcblk0
\end{bashinput}

If \code{cfdisk} asks you to \code{Select a label type}, choose
\code{dos}. This corresponds to traditional partition tables that
DOS/Windows would understand. \code{gpt} partition tables are needed
for disks bigger than 2 TB.

In the \code{cfdisk} interface, delete existing partitions, then
create two partitions with the following properties:

\begin{itemize}
\item First Partition
\begin{itemize}
  \item Size: \code{128MB}
  \item Select \code{primary} partition
  \item Type: \code{W95 FAT32 (LBA)} (\code{c} choice)
  \item Bootable flag enabled
\end{itemize}

\item Second Partition
\begin{itemize}
  \item Size: use the rest of the available space
  \item Select \code{primary} partition
  \item Type: \code{Linux} (\code{83} choice)
\end{itemize}
\end{itemize}

Press \code{Write} when you are done.

To make sure that partition definitions are reloaded on your
workstation, remove the SD card and insert it again.

Now create a FAT32 filesystem on the first partition and an ext4
filesystem on the second one:

\begin{bashinput}
sudo mkfs.vfat -F 32 -n boot /dev/mmcblk0p1
sudo mkfs.ext4 -L rootfs -E nodiscard /dev/mmcblk0p2
\end{bashinput}

You can now make your workstation automatically mount these partitions
by removing the SD card and plugging it back. They should now be
mounted on \code{/media/$USER/boot} and \code{/media/$USER/rootfs}.
